%% Compile with: latexmk -xelatex -interaction=nonstopmode main.tex
\documentclass[11pt,a4paper]{article}

\usepackage[fontsize=10.5pt]{scrextend}
\usepackage{fontspec}
\usepackage{libertinus}
\usepackage{geometry}
\usepackage{titlesec}
\usepackage{enumitem}
\usepackage{xcolor}
\usepackage{microtype}
\usepackage{hyperref}

\geometry{top=1.3cm, bottom=1.3cm, left=2.2cm, right=2.2cm}

\definecolor{linkink}{RGB}{22,52,104}
\definecolor{ruleink}{gray}{0.35}
\definecolor{quietink}{gray}{0.30}

\hypersetup{
  colorlinks=true,
  urlcolor=linkink,
  linkcolor=linkink,
  pdftitle={Chang He (Aaron) -- Curriculum Vitae},
  pdfauthor={Chang He},
  pdfsubject={Curriculum Vitae},
  pdfkeywords={quantum computing, theoretical computer science, algebra}
}

% Section style: letterspaced small caps over a hairline rule
\titleformat{\section}
  {\normalfont\large\scshape}
  {}
  {0em}
  {\addfontfeatures{LetterSpace=6.0}\MakeLowercase}
  [\vspace{-5pt}{\color{ruleink}\rule{\linewidth}{0.4pt}}]
\titlespacing{\section}{0pt}{10pt}{5pt}

\setlist[itemize]{leftmargin=1.4em, topsep=3pt, itemsep=1.5pt, parsep=0pt,
                  label={\raisebox{0.15ex}{\footnotesize\textbullet}}}

\setlength{\parindent}{0pt}
\pagestyle{empty}
\linespread{1.0}
\emergencystretch=1.5em

% \cventry{institution}{place}{role}{dates}
\newcommand{\cventry}[4]{%
  \textbf{#1}\hfill #2\par
  \vspace{1pt}%
  \textit{#3}\hfill \textit{#4}\par}

% \interest{area}{description}
\newcommand{\interest}[2]{%
  \item[\textit{#1}] #2}

\newcommand{\sep}{\hspace{0.6em}{\color{ruleink}\textperiodcentered}\hspace{0.6em}}

\begin{document}

%% ── Header ──────────────────────────────────────────────────────────────────
\begin{center}
  {\huge\scshape\addfontfeatures{LetterSpace=8.0}chang he}%
  \hspace{0.5em}{\Large\itshape (Aaron)}\\[9pt]
  {\small\color{quietink}%
   Undergraduate, \href{https://www.unsw.edu.au}{UNSW Sydney}%
   \sep Sydney, Australia}\\[5pt]
  {\small
   \href{mailto:aaron.he@student.unsw.edu.au}{aaron.he@student.unsw.edu.au}%
   \sep
   \href{https://aaronhnoraa.github.io}{aaronhnoraa.github.io}%
   \sep
   \href{https://github.com/AaronHnoraA}{github.com/AaronHnoraA}%
   \sep
   \href{https://wiki.pwo101.top/wiki}{wiki.pwo101.top}}
\end{center}

%% ── Education ───────────────────────────────────────────────────────────────
\section{Education}

\cventry
  {\href{https://www.unsw.edu.au}{University of New South Wales (UNSW)}}
  {Sydney, Australia}
  {Bachelor of Advanced Computer Science (Honours)
   (\href{https://www.handbook.unsw.edu.au/undergraduate/programs/2026/3779}{UNSW 3779})}
  {Term 3, 2024 -- Present}

\begin{itemize}
  \item Expected graduation: Term 3, 2028. Transferred from Bachelor of Science in Mathematics
        (\href{https://www.handbook.unsw.edu.au/undergraduate/programs/2026/3956}{UNSW 3956}).
  \item Academic standing: WAM 83.5; admitted to the UNSW
        \href{https://www.unsw.edu.au/science/student-life-resources/student-opportunities/talented-students-program}{Talented Students Program}
        (TSP) in 2025.
\end{itemize}

%% ── Research Interests ──────────────────────────────────────────────────────
\section{Research Interests}

\begin{description}[leftmargin=0pt, labelsep=0.5em, topsep=0pt, itemsep=3pt,
                    parsep=0pt, font=\normalfont]
  \interest{Quantum computing.}
    {Quantum algorithms and quantum simulation.}
  \interest{Theoretical computer science.}
    {Computational complexity and algebraic algorithms: tensor, group, and graph
     isomorphism; linear code equivalence (LCE); reductions among equivalence
     problems.}
  \interest{Algebra.}
    {Multilinear algebra and tensor spaces, group actions, canonical forms and
     invariants.}
\end{description}

%% ── Research Experience ─────────────────────────────────────────────────────
\section{Research Experience}

\cventry
  {Research, \href{https://www.uts.edu.au}{University of Technology Sydney (UTS)}}
  {Sydney, Australia}
  {Supervised by \href{https://sites.google.com/site/jimmyqiao86/}{Youming Qiao}}
  {October 2026 -- Present}

\begin{itemize}
  \item Starting work on the linear code equivalence (LCE) and matrix code
        equivalence (MCE) problems.
\end{itemize}

\vspace{4pt}

\cventry
  {Research, \href{https://www.uts.edu.au}{University of Technology Sydney (UTS)}}
  {Sydney, Australia}
  {Supervised by \href{https://sites.google.com/site/jimmyqiao86/}{Youming Qiao}
   and \href{https://sites.google.com/site/melderau/}{Murray Elder}}
  {Late 2025 -- September 2026}

\begin{itemize}
  \item Studied isomorphism and automorphism problems for tensors, groups, algebras, and polynomials.
  \item Proved that an $r$-uniform hypergraph is connected if and only if its
        associated tensor is indecomposable, linking graph structure to tensor
        structure; contributed to the manuscript below.
\end{itemize}

\vspace{4pt}

\cventry
  {Independent Study, \href{https://www.uts.edu.au}{University of Technology Sydney (UTS)}}
  {Sydney, Australia}
  {Supervised by \href{https://sites.google.com/site/jimmyqiao86/}{Youming Qiao}}
  {Term 2 -- Term 3, 2025}

\begin{itemize}
  \item Supervised reading in quantum computing: Nielsen and Chuang,
        \textit{Quantum Computation and Quantum Information}, and other standard texts.
\end{itemize}

\vspace{4pt}

\cventry
  {\href{https://www.unsw.edu.au/science/student-life-resources/student-opportunities/talented-students-program}{Talented Students Program}, UNSW}
  {Sydney, Australia}
  {Supervised by \href{https://research.unsw.edu.au/people/dr-quoc-thong-le-gia}{Quoc Le Gia}}
  {Term 1, 2025}

\begin{itemize}
  \item Supervised reading of selected literature in large language models; practice in
        extracting technical assumptions and presenting mathematical and algorithmic ideas.
\end{itemize}

%% ── Manuscripts ─────────────────────────────────────────────────────────────
\section{Manuscripts}

\begin{itemize}
  \item Murray Elder, \textbf{Chang He}, Luke Mathieson, Euan Mendoza, and Youming Qiao.
        \textit{On isomorphism and automorphism problems for tensors, groups,
        algebras, and polynomials.}
        Submitted to ITCS 2027, under review (authors in alphabetical order).
\end{itemize}

%% ── Technical Skills ────────────────────────────────────────────────────────
\section{Technical Skills}

\begin{description}[leftmargin=0pt, labelsep=0.5em, topsep=0pt, itemsep=2pt,
                    parsep=0pt, font=\normalfont]
  \interest{Programming.}{Python, C/C++, JavaScript/TypeScript; Git, LaTeX, Unix shell.}
  \interest{Quantum and AI.}{Qiskit (working knowledge); building and using LLM agents.}
\end{description}

%% ── Academic Activities ─────────────────────────────────────────────────────
\section{Academic Activities}

\begin{itemize}
  \item Member, Sydney Algorithms and Computing Theory (SACT) group, University of Sydney.
  \item \href{https://www.maths.usyd.edu.au/s/scnitm/pieter-QuantumComputingReadingGr}{MathQuEST
        Reading Group on Quantum Computing}, University of Sydney, 2026.
  \item \href{https://sites.google.com/view/sydney-tcs-winter-school-2026/home}{Sydney
        TCS Winter School}: Interactive Proofs and the PCP Theorem, University of Sydney, July 2026.
\end{itemize}

\end{document}
